%====================
% EDUCATION
%====================

\section{Education}

\subsection{{Bachelor of Science in Electrical Engineering \hfill August 2015 – May 2019}}
\subtext{University of Illinois at Urbana–Champaign \hfill Champaign-Urbana, IL}
\begin{zitemize}
    \item Relevant coursework: Computer Programming, Electromagnetics, Semiconductor Physics, Analog and Digital Circuits.
    \item Internships: GE Global Research (Summer 2016 \& 2017), Amazon (Summer 2018).
    \item Exchange Semester: Technical University of Denmark (Autumn 2017).
\end{zitemize}
\vspace{6pt}

%====================
% TECHNICAL EXPERIENCE
%====================

\section{Technical Experience}

\subsection{{Research Coordinator (Photonics) \hfill May 2024 – May 2025}}
\subtext{University of Washington, College of the Environment, Department of Earth \& Space Sciences \hfill Seattle, WA}
\begin{zitemize}
    \item Operated, evaluated, and deployed fiber-optic interrogator systems (Sintella Onyx, ASN OptoDAS, FEBUS G1-C) for DAS/DTS/DSS sensing modes across seismic, ocean acoustic, and glacier/ice-ocean dynamics applications.
    \item Conducted field deployments across six sites including offshore fiber installations in Homer AK, San Juan Islands WA, and Morro Bay CA, and seismic monitoring deployments at Mt. Rainier and Alpental Ski Resort WA; handled all pre-deployment acceptance testing, independent equipment transport, live-buffer drive swaps, and ferry and vehicle logistics.
    \item Maintained daily operational oversight of concurrent DAS deployments across domestic and international field sites, managing data acquisition systems, remote monitoring infrastructure, and storage logistics across petabyte-scale datasets.
    \item Configured remote network access for deployed interrogator systems, coordinated RAID storage architecture decisions with university IT, and ensured federal data handling compliance including Navy data scrubbing protocols for oceanographic datasets.
    \item Coordinated international equipment shipments and customs carnets for Antarctic and Greenland deployments; managed procurement and lab inventory; facilitated dataset publication through DOI creation with UW Libraries; collaborated closely with USGS, EarthScope, and PNSN; contributed to and credited in Shi et al. 2025, Seismological Research Letters.
    \item Actively participated in research group colloquia and engaged in future planning discussions; contributed to signal analysis including data cross-correlation and live characterization datasets.
\end{zitemize}
\vspace{6pt}
\subsection{{Hardware Engineer II (Signal Integrity) \hfill September 2021 – September 2023}}
\subtext{Microsoft, Cloud AI Hardware and Advanced Signal Engineering \hfill Seattle, WA}
\begin{zitemize}
    \item Performed passive and active signal integrity characterization of PCIe and DDR5 traces across FPGA die, custom silicon packaging, connectors, and board-level interconnects using differential probing, VNA S-parameter analysis, and time-domain eye-diagram methods.
    \item Identified component-level signal integrity defects through comparative validation of boards across manufacturers and production origins, enabling root cause isolation and cost savings in hardware qualification.
    \item Supported stackup design decisions alongside motherboard EEs, contributing to material selection and layer thickness tradeoffs that determined high-speed signal propagation behavior.
    \item Built end-to-end simulation models of signal paths from silicon to GPU and PCIe endpoints using ANSYS HFSS and Keysight ADS, analyzing microbend and geometry effects on transmission performance.
\end{zitemize}
\vspace{6pt}
\subsection{{Hardware Design Engineer (Satellite Avionics) \hfill April 2020 – September 2021}}
\subtext{Amazon, Project Kuiper \hfill Seattle, WA}
\begin{zitemize}
    \item Owned end-to-end hardware development for magnetometer and sun sensor subsystems during early-phase pre-PDR program development (prior to launch provider selection and during initial lab standup) including electrical, mechanical, and functional requirements definition, FOV analysis, and ConOps documentation.
    \item Led component selection, schematic design, and PCB layout for avionics sensor prototypes; coordinated with contract manufacturers on assembly requirements and parts library verification.
    \item Performed sensor bring-up, functional validation, and photodiode characterization using test tools such as sun emulator, thermal chamber, vibration test bed, and Helmholtz coil; conducted field validation with standalone prototype hardware.
    \item Designed and populated custom test boards for heavy-ion single-event effects (SEE) characterization at the Lawrence Berkeley National Laboratory 88-Inch Cyclotron; prepared delidded commercial components, coordinated beamline exposure parameters, and performed pre- and post-irradiation functional testing of semiconductor devices; contributed to Kuiper's internal COTS component radiation qualification library.
    \item Supported PDR and CDR design reviews and contributed to cross-team integration reviews for bus avionics sensor suite.
    \item Maintained comprehensive engineering documentation across requirements, design decisions, test results, and manufacturing handoff throughout full hardware development lifecycle.
\end{zitemize}
\vspace{6pt}

%====================
% SKILLS & CERTIFICATIONS
%====================

\section{Skills and Certifications}

\renewcommand{\arraystretch}{1.2}

\hyphenpenalty=10000
\exhyphenpenalty=10000

\begin{tabularx}{\linewidth}{p{4.8cm} X}

\skills{Languages} & English (native), French (A2), Spanish (A2) \\

\skills{Instrumentation} & \textit{Photonics}\normalfont{: DAS/DTS/DSS, Optical Connectors \& Fiber Components, Fusion Splicing}
 \\

  & \textit{Interrogators}\normalfont{: Sintella Onyx, ASN OptoDAS, FEBUS G1-C}
 \\

  & \textit{Electronic Test~\&~Measurement}\normalfont{: VNA, S-parameters, Eye-diagram Analysis, Oscilloscope, Passive \& Active Probing, Photodetector Characterization, EMI/EMC Validation, Thermal/Mechanical Environmental Qualification}
 \\

\skills{Simulation~\&~Design} & \textit{Hardware Design}\normalfont{: Schematic, PCB Layout, Board Bring-up, Stackup Design, PCB Fabrication \& Assembly Processes}
 \\

  & \textit{Simulation~\&~Analysis}\normalfont{: ANSYS HFSS, Keysight ADS, Signal Integrity Analysis, Channel Simulation, Noise Characterization}
 \\

\skills{Programming~\&~Development} & \textit{Languages~\&~Tools}\normalfont{: Python, MATLAB, C/C++, JavaScript, LaTeX, Git, Linux/Bash, Docker, Jupyter}
 \\

  & \textit{Infrastructure}\normalfont{: RAID Storage Architecture, Data Acquisition Pipelines, Remote System Monitoring}
 \\

\skills{Domain Knowledge} & \textit{Engineering}\normalfont{: Electromagnetics, Signal Integrity, SEE Radiation Effects Testing, COTS Component Qualification}
 \\

  & \textit{Research}\normalfont{: Geophysical Sensing, Ocean Acoustics, Cryospheric Instrumentation, Seismic Monitoring}
 \\

\skills{Field~\&~Operations} & \textit{Logistics}\normalfont{: International Shipping, Customs \& Carnets, Lab Budget \& Procurement, Inventory Management, Project Tracking \& Milestone Management}
 \\

  & \textit{Research Operations}\normalfont{: Federal Data Compliance, Petabyte-scale Data Management, Grant Management \& Documentation, Technical Reporting, Multi-institution \& International Research Collaboration}
 \\

\end{tabularx}

\hyphenpenalty=10000
\exhyphenpenalty=10000

\begin{tabularx}{\linewidth}{p{4.8cm} X}
\skills{Field Activities} & American Mountain Guides Association (AMGA) Professional Member, American Alpine Club (AAC) Member, Avalanche Rescue Training, Wilderness First Responder (WFR) with AED and CPR Certifications, OSHA 10 Training \\
\skills{Workshops} & DesignCon (2022), BOAT Ocean Acoustics Workshop (2025) \\
\end{tabularx}